% Source: https://artofproblemsolving.com/community/c3h2412052_polynomials_urk
% Stable identity: aops-topic-2412052-q01
\textbf{(a)} If $\dfrac pq$ (in lowest terms, $q>0$) is a root, then $p$ divides the constant term $60$ and $q$ divides the leading coefficient $12$. The divisors of $60$ are $1,2,3,4,5,6,10,12,15,20,30,60$ and those of $12$ are $1,2,3,4,6,12$. The possible rational roots are $\pm r$, where $r$ is one of the $30$ different fractions
\[\dfrac1{12},\dfrac16,\dfrac14,\dfrac13,\dfrac5{12},\dfrac12,\dfrac23,\dfrac34,\dfrac56,1,\dfrac54,\dfrac43,\dfrac32,\dfrac53,2,\dfrac52,3,\dfrac{10}3,\dfrac{15}4,4,5,6,\dfrac{20}3,\dfrac{15}2,10,12,15,20,30,60,\]
listed in increasing order: $60$ candidates in all.

\textbf{(b)} \emph{Positive roots.} For $0<x\le1$, $f(x)\le12-x-60<0$. For $x\ge1$, $f'(x)=48x^3-1>0$, so $f$ is increasing there. Since
\[f\left(\dfrac32\right)=\dfrac{243}{4}-\dfrac32-60=-\dfrac34<0,\qquad f\left(\dfrac53\right)=\dfrac{2500}{27}-\dfrac53-60=\dfrac{835}{27}>0,\]
$f$ has exactly one positive root, and it lies strictly between $\dfrac32$ and $\dfrac53$, which are consecutive numbers in the list. So it is not a candidate.

\emph{Negative roots.} Write $x=-t$ with $t>0$: $f(-t)=12t^4+t-60$, which is increasing in $t$, with
\[f\left(-\dfrac43\right)=\dfrac{1024}{27}+\dfrac43-60=-\dfrac{560}{27}<0,\qquad f\left(-\dfrac32\right)=\dfrac{243}{4}+\dfrac32-60=\dfrac94>0.\]
So the only negative root lies strictly between $-\dfrac32$ and $-\dfrac43$, and again no candidate lies there ($\dfrac43$ and $\dfrac32$ are consecutive in the list).

Hence $\boxed{f\text{ has no rational roots}}$: none of the $60$ candidates is a root. (Its real roots are about $1.505$ and $-1.486$; the other two roots are not real.)
